\section{Diagnóstico y descripción de la realidad}
\label{sec:diagnostico}

El MÓDULO TECNOLÓGICO PRODUCTIVO SAN JULIÁN, ubicado en el municipio de San Julián del departamento de Santa Cruz, constituye la institución formativa de referencia del Bachillerato Técnico Humanístico (BTH) para los estudiantes del nivel secundario procedentes de las diversas unidades educativas del distrito escolar. Cada inicio de gestión académica, la demanda de cupos para ingresar a cualquiera de sus siete especialidades técnicas (Construcción Civil, Contaduría Pública, Electrónica, Gastronomía, Mecánica Automotriz, Sistemas Informáticos y Textiles y Confecciones) sobrepasa ampliamente la capacidad instalada de aulas y talleres institucionales, generando una intensa competencia por las vacantes habilitadas.

El diagnóstico de la realidad operativa de la institución consideró el análisis retrospectivo de las gestiones 2022 a 2025 y un relevamiento empírico de campo realizado durante el primer bimestre de la gestión 2026, combinando encuestas aplicadas a estudiantes postulantes y entrevistas en profundidad con directivos, docentes y personal administrativo. En una primera etapa, correspondiente a las gestiones presenciales tradicionales de 2022 a 2024, el registro se realizaba de forma manuscrita en ventanillas físicas. El temor generalizado a no alcanzar un cupo en la carrera deseada forzaba a los estudiantes y a sus familias a acampar y formar extensas filas a la intemperie en los alrededores del módulo con varios días de anticipación, sufriendo desgaste físico, exposición a factores climáticos adversos y abandono temporal de sus actividades laborales agrícolas y comerciales.

Con el fin de superar estas problemáticas y dar un paso hacia la digitalización, en la gestión 2025 se habilitó una primera versión del sistema web de inscripción. Sin embargo, dicha plataforma evidenció una limitación funcional crítica: carecía en absoluto de un mecanismo de recuperación autónoma de contraseñas o de consulta de usuario asignado. Ante el olvido recurrente de credenciales de acceso, cientos de postulantes debieron concurrir de manera masiva y urgente a las instalaciones físicas del módulo para solicitar el restablecimiento manual de sus cuentas ante el personal administrativo. Durante el tiempo consumido en estas largas filas de reseteo presencial, los cupos disponibles en las carreras de mayor preferencia se agotaron, obligando a numerosos estudiantes a resignarse a matricularse en especialidades técnicas que no representaban su verdadera vocación profesional.

De manera simultánea, el diagnóstico reveló una profunda ruptura en la continuidad del ciclo administrativo: mientras que el registro web preliminar se llevaba a cabo al cierre de la gestión previa como una reserva de vacante, la administración no contaba con ningún sistema informático para formalizar y consolidar la matriculación definitiva en la gestión académica siguiente. En el régimen institucional del BTH, la matriculación efectiva exige que el estudiante acredite documentalmente sus libretas de calificaciones de 1ro a 3ro de secundaria, la cédula de identidad de sus tutores o apoderados, la fotocopia de su carnet de identidad y su certificado de nacimiento. La inexistencia de un módulo de software para recepcionar, verificar y registrar digitalmente estos requisitos generó un descontrol en el archivo de carpetas físicas, expedientes incompletos y una excesiva sobrecarga de trabajo manual para la secretaría académica.

Finalmente, el levantamiento de información constató la ausencia total de un mecanismo automatizado para la asignación y organización de los estudiantes en cursos y paralelos específicos. Tras consolidar la inscripción en una especialidad técnica, la administración del módulo enfrenta la tarea de distribuir a los matriculados en sus respectivos cursos y paralelos (tales como 4to A, 4to B u otros grados) tomando como criterio rector la unidad educativa de procedencia a la que pertenece cada estudiante. Al ejecutarse este procedimiento de forma empírica y manual mediante planillas externas, se originaron constantes errores de transcripción, sobrepoblación en ciertos talleres y graves conflictos de coordinación horaria entre el módulo tecnológico y las distintas unidades educativas del distrito municipal.
